% GENERATED FILE. Edit canonical lessons, then run npm run generate:books.
% canonical-source-hash: fnv1a64:fec0637a9b0335a3
% canonical-lessons: MR-C312-akasvani, MR-C312-durdarsan, MR-C312-vahini, MR-C312-nivedak, MR-C312-prasarit-karne

\chapter{On the Air}
\label{ch:mr-on-the-air}
\hlchaptermodality{312}

\begin{chapteropening}
\textbf{By the end of this chapter:} \emph{I can talk about radio, television and broadcasts.}

Complete the last lesson of chapter 312: I can talk about radio, television and broadcasts.
\end{chapteropening}

\section[\texorpdfstring{\mr{आकाशवाणी} (ākāśvāṇī)}{आकाशवाणी (ākāśvāṇī)}]{\mr{आकाशवाणी} (ākāśvāṇī) — radio; the radio service}

\label{lesson:MR-C312-akasvani}



\begin{quote}
Before the new one: say the Marathi for an editorial, then the Marathi for to publish.
\end{quote}

\subsection*{You'll want to know: \mr{आकाशवाणी}}
\textbf{\mr{आकाशवाणी}} — \emph{ākāśvāṇī} — \textquotedblleft{}radio; the radio service\textquotedblright{}.

\begin{quote}
\mr{आजी} \mr{रोज} \mr{सकाळी} \mr{आकाशवाणीवर} \mr{गाणी} \mr{ऐकते}.
\end{quote}

\begin{quote}
\emph{ājī roj sakāḷī ākāśvāṇīvar gāṇī aikte.}
\end{quote}

\begin{quote}
Grandma listens to songs on the radio every morning.
\end{quote}

Feminine.

\mr{आकाशवाणीवर}, on the radio. \mr{रेडिओ} is also said.

\begin{cousinweb}
\textbf{\mr{आकाशवाणी}} (\emph{ākāśvāṇī}) is a tatsama, taken from Sanskrit \emph{ākāśavāṇī}, 'a voice from the air or from heaven' (Monier-Williams), a fitting name for the radio.
\end{cousinweb}

\begin{grammarlens}[title={What you've built}]
One more word for telling what happened.
\end{grammarlens}

\subsection*{Your turn}
\begin{itemize}
\raggedright
  \item \emph{Say it:} \emph{ākāśvāṇī}
  \item \emph{Say it:} \emph{ākāśvāṇī}, once more
  \item \emph{Say it:} \emph{ājī roj sakāḷī ākāśvāṇīvar gāṇī aikte.} — Grandma listens to songs on the radio every morning.
  \item \emph{Recall:} say the Marathi for an editorial, then the Marathi for to publish, then say \emph{ākāśvāṇī} again
\end{itemize}

\begin{tcolorbox}[breakable,colback=teal!4,colframe=teal!35!black,title={Before you move on}]
What does \mr{आकाशवाणी} mean? (\textquotedblleft{}Radio; the radio service\textquotedblright{}.) Say it once more.
\end{tcolorbox}
\section[\texorpdfstring{\mr{दूरदर्शन} (dūrdarśan)}{दूरदर्शन (dūrdarśan)}]{\mr{दूरदर्शन} (dūrdarśan) — television}

\label{lesson:MR-C312-durdarsan}



\begin{quote}
Before the new one: say the Marathi for to publish, then the Marathi for radio.
\end{quote}

\subsection*{You'll want to know: \mr{दूरदर्शन}}
\textbf{\mr{दूरदर्शन}} — \emph{dūrdarśan} — \textquotedblleft{}television\textquotedblright{}.

\begin{quote}
\mr{रात्री} \mr{आम्ही} \mr{दूरदर्शनवर} \mr{बातम्या} \mr{पाहतो}.
\end{quote}

\begin{quote}
\emph{rātrī āmhī dūrdarśanvar bātamyā pāhto.}
\end{quote}

\begin{quote}
At night we watch the news on television.
\end{quote}

Neuter.

Also the name of the national channel; in speech \mr{टीव्ही}.

\begin{cousinweb}
\textbf{\mr{दूरदर्शन}} (\emph{dūrdarśan}) is a modern coinage built from Sanskrit \emph{dūra}, 'far', and \emph{darśana}, 'seeing' (Monier-Williams): seeing far.
\end{cousinweb}

\begin{grammarlens}[title={What you've built}]
One more word for telling what happened.
\end{grammarlens}

\subsection*{Your turn}
\begin{itemize}
\raggedright
  \item \emph{Say it:} \emph{dūrdarśan}
  \item \emph{Say it:} \emph{dūrdarśan}, once more
  \item \emph{Say it:} \emph{rātrī āmhī dūrdarśanvar bātamyā pāhto.} — At night we watch the news on television.
  \item \emph{Recall:} say the Marathi for to publish, then the Marathi for radio, then say \emph{dūrdarśan} again
\end{itemize}

\begin{tcolorbox}[breakable,colback=teal!4,colframe=teal!35!black,title={Before you move on}]
What does \mr{दूरदर्शन} mean? (\textquotedblleft{}Television\textquotedblright{}.) Say it once more.
\end{tcolorbox}
\section[\texorpdfstring{\mr{वाहिनी} (vāhinī)}{वाहिनी (vāhinī)}]{\mr{वाहिनी} (vāhinī) — a channel}

\label{lesson:MR-C312-vahini}



\begin{quote}
Before the new one: say the Marathi for radio, then the Marathi for television.
\end{quote}

\subsection*{You'll want to know: \mr{वाहिनी}}
\textbf{\mr{वाहिनी}} — \emph{vāhinī} — \textquotedblleft{}a channel\textquotedblright{}.

\begin{quote}
\mr{ही} \mr{वाहिनी} \mr{दिवसभर} \mr{फक्त} \mr{बातम्या} \mr{दाखवते}.
\end{quote}

\begin{quote}
\emph{hī vāhinī divasbhar phakta bātamyā dākhavte.}
\end{quote}

\begin{quote}
This channel shows only news all day.
\end{quote}

Feminine.

\mr{वृत्तवाहिनी}, a news channel.

\begin{cousinweb}
\textbf{\mr{वाहिनी}} (\emph{vāhinī}) is a tatsama, taken from Sanskrit \emph{vāhinī}, 'an army, a host' (Monier-Williams), a word built on \emph{vah}, 'to carry'; the sense 'a channel' is modern.
\end{cousinweb}

\begin{grammarlens}[title={What you've built}]
One more word for telling what happened.
\end{grammarlens}

\subsection*{Your turn}
\begin{itemize}
\raggedright
  \item \emph{Say it:} \emph{vāhinī}
  \item \emph{Say it:} \emph{vāhinī}, once more
  \item \emph{Say it:} \emph{hī vāhinī divasbhar phakta bātamyā dākhavte.} — This channel shows only news all day.
  \item \emph{Recall:} say the Marathi for radio, then the Marathi for television, then say \emph{vāhinī} again
\end{itemize}

\begin{tcolorbox}[breakable,colback=teal!4,colframe=teal!35!black,title={Before you move on}]
What does \mr{वाहिनी} mean? (\textquotedblleft{}A channel\textquotedblright{}.) Say it once more.
\end{tcolorbox}
\section[\texorpdfstring{\mr{निवेदक} (nivedak)}{निवेदक (nivedak)}]{\mr{निवेदक} (nivedak) — a newsreader; a presenter}

\label{lesson:MR-C312-nivedak}



\begin{quote}
Before the new one: say the Marathi for television, then the Marathi for a channel.
\end{quote}

\subsection*{You'll want to know: \mr{निवेदक}}
\textbf{\mr{निवेदक}} — \emph{nivedak} — \textquotedblleft{}a newsreader; a presenter\textquotedblright{}.

\begin{quote}
\mr{निवेदकाने} \mr{शांत} \mr{आवाजात} \mr{बातमी} \mr{वाचली}.
\end{quote}

\begin{quote}
\emph{nivedakāne śānt āvājāt bātmī vāclī.}
\end{quote}

\begin{quote}
The newsreader read the news in a calm voice.
\end{quote}

Masculine.

A woman presenter is \mr{निवेदिका}.

\begin{cousinweb}
\textbf{\mr{निवेदक}} (\emph{nivedak}) is a tatsama, taken from Sanskrit \emph{nivedaka}, 'communicating, relating' (Monier-Williams).
\end{cousinweb}

\begin{grammarlens}[title={What you've built}]
One more word for telling what happened.
\end{grammarlens}

\subsection*{Your turn}
\begin{itemize}
\raggedright
  \item \emph{Say it:} \emph{nivedak}
  \item \emph{Say it:} \emph{nivedak}, once more
  \item \emph{Say it:} \emph{nivedakāne śānt āvājāt bātmī vāclī.} — The newsreader read the news in a calm voice.
  \item \emph{Recall:} say the Marathi for television, then the Marathi for a channel, then say \emph{nivedak} again
\end{itemize}

\begin{tcolorbox}[breakable,colback=teal!4,colframe=teal!35!black,title={Before you move on}]
What does \mr{निवेदक} mean? (\textquotedblleft{}A newsreader; a presenter\textquotedblright{}.) Say it once more.
\end{tcolorbox}
\section[\texorpdfstring{\mr{प्रसारित} \mr{करणे} (prasārit karṇe)}{प्रसारित करणे (prasārit karṇe)}]{\mr{प्रसारित} \mr{करणे} (prasārit karṇe) — to broadcast}

\label{lesson:MR-C312-prasarit-karne}



\begin{quote}
Before the new one: say the Marathi for a channel, then the Marathi for a newsreader.
\end{quote}

\subsection*{You'll want to know: \mr{प्रसारित} \mr{करणे}}
\textbf{\mr{प्रसारित} \mr{करणे}} — \emph{prasārit karṇe} — \textquotedblleft{}to broadcast\textquotedblright{}.

\begin{quote}
\mr{वाहिनीने} \mr{ती} \mr{मुलाखत} \mr{संध्याकाळी} \mr{प्रसारित} \mr{केली}.
\end{quote}

\begin{quote}
\emph{vāhinīne tī mulākhat sandhyākāḷī prasārit kelī.}
\end{quote}

\begin{quote}
The channel broadcast that interview in the evening.
\end{quote}

Past, with -\mr{ने} (he …): \mr{त्याने} \mr{प्रसारित} \mr{केलं}, \emph{tyāne prasārit kelaṁ}.

What is broadcast agrees: \mr{मुलाखत} \mr{प्रसारित} \mr{केली}, \mr{कार्यक्रम} \mr{प्रसारित} \mr{केला}.

\begin{cousinweb}
\textbf{\mr{प्रसारित} \mr{करणे}} (\emph{prasārit karṇe}) joins \emph{prasārit}, a tatsama taken from Sanskrit \emph{prasārita}, 'spread out, stretched forth' (Monier-Williams), to \emph{karṇe}, 'to do'.
\end{cousinweb}

\begin{grammarlens}[title={What you've built}]
One more word for telling what happened.
\end{grammarlens}

\subsection*{Your turn}
\begin{itemize}
\raggedright
  \item \emph{Say it:} \emph{prasārit karṇe}
  \item \emph{Say it:} \emph{prasārit karṇe}, once more
  \item \emph{Say it:} \emph{vāhinīne tī mulākhat sandhyākāḷī prasārit kelī.} — The channel broadcast that interview in the evening.
  \item \emph{Recall:} say the Marathi for a channel, then the Marathi for a newsreader, then say \emph{prasārit karṇe} again
\end{itemize}

\begin{tcolorbox}[breakable,colback=teal!4,colframe=teal!35!black,title={Before you move on}]
What does \mr{प्रसारित} \mr{करणे} mean? (\textquotedblleft{}To broadcast\textquotedblright{}.) Say it once more.
\end{tcolorbox}
